\documentclass[10pt]{article}

% graphicx package, useful for including eps, pdf, and raster graphics
\usepackage{graphicx}
\DeclareGraphicsExtensions{.png,.png,.jpg,.pdf}

% basic packages
\usepackage{fancyhdr}
\usepackage{color}
\usepackage{parskip}
\usepackage{float}
\usepackage{microtype}
\usepackage{url}
\usepackage{hyperref}
\usepackage{booktabs}
\usepackage{makecell}
\usepackage{multirow}
\usepackage{pbox}
\usepackage{enumitem}

% page layout
\usepackage[top=1in, bottom=1in, left=0.75in, right=0.75in, includefoot, heightrounded]{geometry}

% Useful for seeing the layout as you tweak it.
%\usepackage{showframe}

% fonts
\usepackage{fontspec}

\usepackage{polyglossia}
\setmainlanguage{vietnamese}
\setotherlanguage{english}

\setmainfont{DroidSans}[
  Extension   = .ttf,
  UprightFont = *,
  BoldFont    = *-Bold ]
\setmonofont{DroidSansMono}[
  Extension   = .ttf,
  UprightFont = * ]

\hypersetup{colorlinks,urlcolor=blue}
\urlstyle{same}

\setlength{\headheight}{1.1in}
\pagestyle{fancy}
\fancyhf{}

\lhead{\includegraphics[height=1in,clip]{logo}}
\rhead{\pbox[b][1in][c]{\textwidth}{
    \small
    \textbf{Xét Nghiệm Được Thực Hiện Bởi:}\\
    Northwest Genomics Center\\
    University of Washington School of Medicine\\
    Genome Sciences Box 355065\\
    3720 15th Ave NE\\
    Seattle, WA 98195-5065}}

\newcommand{\PageLine}{\rule{\textwidth}{0.25mm}}

\begin{document}

\begin{description}[font=\normalfont,align=left,labelwidth=10em]
\item [\textbf{\VAR{pat_name|e}}]
\item [\textbf{Ngày Sinh:}] \VAR{birth_date|e}
\item [\textbf{Mẫu Bệnh Phẩm:}] Giữa mũi, tự lấy
\item [\textbf{Mã Vạch Lấy Mẫu:}] \VAR{qrcode|e}
\item [\textbf{Ngày Lấy Mẫu:}] \VAR{collect_ts|e}
\item [\textbf{Ngày Báo Cáo:}]
  %- if status_code == "never-tested"
  không bao giờ xét nghiệm
  %- else
  \VAR{result_ts|e}
  %- endif
\end{description}

\PageLine

\begin{center}
\Large
\textbf{KẾT QUẢ NGHIÊN CỨU - Khuyến khích quý vị chia sẻ thông tin này với nhà
  cung cấp dịch vụ chăm sóc sức khỏe của mình}
\end{center}

\bigskip

\large \underline{\textbf{Giới Thiệu về Xét Nghiệm}}

\begin{itemize}
\item

  COVID-19 là \textbf{căn bệnh} do \textbf{vi-rút} SARS-CoV-2 gây ra. Xét nghiệm
  mang tính nghiên cứu này tìm vi-rút trong mẫu bệnh phẩm được thu thập từ mũi
  của quý vị để xác định liệu quý vị có mắc COVID-19 hay không. Thuật ngữ kỹ
  thuật cho loại xét nghiệm này là ``phản ứng tổng hợp chuỗi phiên mã ngược''
  hay ngắn gọn là ``RT-PCR''.

\end{itemize}

\bigskip

\large \underline{\textbf{Nhắc Nhở Quan Trọng}}

Trước khi xem kết quả xét nghiệm mang tính nghiên cứu của quý vị, quý vị nên xem
xét những thông tin sau:

\begin{itemize}
\item

  \textbf{Loại mẫu} mà quý vị thu thập và \textbf{xét nghiệm trong phòng thí
  nghiệm này} là một phần trong \textbf{nghiên cứu} và báo cáo này mô tả
  \textbf{kết quả xét nghiệm mang tính nghiên cứu} của quý vị.

\item

  Quý vị đã tự lấy mẫu và dán nhãn cho bản thân mình (hoặc làm cho ai đó), sau
  đó gửi lại cho phòng thí nghiệm qua dịch vụ giao hàng. Mẫu không được thu
  thập, dán nhãn hoặc giao nhận bởi một chuyên gia chăm sóc sức khỏe được cấp
  phép như y tá hoặc bác sĩ. Điều này có thể ảnh hưởng đến kết quả xét nghiệm
  của quý vị và một nhà cung cấp dịch vụ chăm sóc sức khỏe có thể muốn xét
  nghiệm lại cho quý vị.

\item

  Việc trở thành một phần trong \textbf{nghiên cứu} này không thay thế cho các
  chăm sóc y tế thông thường của quý vị. Nếu quý vị cảm thấy ốm nặng hoặc lo
  ngại về kết quả xét nghiệm của mình, quý vị nên liên hệ với một bác sĩ hoặc
  nhà cung cấp dịch vụ chăm sóc sức khỏe.

\end{itemize}

\bigskip
\PageLine

\large \underline{\textbf{Kết Quả Xét Nghiệm Mang Tính Nghiên Cứu SARS-CoV-2
  (COVID-19) Của Quý Vị}}

%- if status_code == "negative"
Không Phát Hiện Thấy Vi-rút (Âm Tính)\\
%- elif status_code == "positive"
Phát Hiện Thấy Vi-rút (Dương Tính)\\
%- elif status_code == "inconclusive"
Không Xác Định\\
%- elif status_code == "never-tested"
Không Thể Xét Nghiệm\\
%- else
\VAR{0/0} %# poor man's raise
%- endif

\PageLine
\bigskip

\large \underline{\textbf{Hiểu Kết Quả Của Quý Vị}}

Bảng này cho biết tất cả các kết quả mà xét nghiệm của quý vị có thể cho thấy và
mô tả ý nghĩa của từng loại kết quả. Quý vị có thể in hoặc gửi email
\textbf{kết quả xét nghiệm mang tính nghiên cứu này} cho bác sĩ hoặc nhà cung
cấp dịch vụ chăm sóc sức khỏe của quý vị.

Bất kể kết quả của quý vị ra sao, quý vị nên ở nhà và tránh xa những người khác
nếu quý vị cảm thấy ốm. \textbf{Nếu các triệu chứng của quý vị tiếp tục, quý vị
cảm thấy ốm nặng hoặc lo lắng về bệnh tình của mình hay có câu hỏi về kết quả
xét nghiệm mang tính nghiên cứu của quý vị, thì quý vị nên liên hệ với một bác
sĩ hoặc nhà cung cấp dịch vụ chăm sóc sức khỏe.}

\begin{itemize}


\item

  \textbf{Nếu xét nghiệm của quý vị là ``Không Phát Hiện Thấy (Âm Tính)''}, điều
  đó có nghĩa là phòng thí nghiệm của chúng tôi không tìm thấy vi-rút gây ra
  COVID-19 trong mẫu mà quý vị gửi cho chúng tôi. Không có xét nghiệm nào hoàn
  hảo và vẫn có khả năng nhỏ là quý vị có thể nhiễm vi-rút này mặc dù kết quả
  xét nghiệm của quý vị là âm tính.

\item

  \textbf{Nếu xét nghiệm của quý vị là ``Phát Hiện Thấy (Dương Tính)''}, điều đó
  có nghĩa là phòng thí nghiệm của chúng tôi tìm thấy vi-rút COVID-19 trong mẫu
  mà quý vị gửi cho chúng tôi. Quý vị nên ở nhà và tránh xa những người khác,
  đặc biệt nếu quý vị cảm thấy ốm.

  Xét nghiệm của quý vị sẽ được chia sẻ với Tiểu Bang Washington và các cơ quan
  y tế địa phương. Quý vị \textbf{không} cần tự báo cáo những kết quả đó. Cơ
  quan y tế có thể gọi cho quý vị để hỏi về bệnh tình của quý vị và những người
  ở gần quý vị trong thời gian quý vị bị ốm.

\item

  \textbf{Nếu xét nghiệm của quý vị là ``Không Xác Định''}, điều đó có nghĩa là
  phòng thí nghiệm của chúng tôi không chắc chắn về kết quả xét nghiệm của quý
  vị. Không có xét nghiệm nào hoàn hảo và kết quả không xác định xảy ra khi
  chúng tôi tìm thấy một số lượng nhỏ SARS-CoV-2 trong mẫu của quý vị, nhưng
  không đủ để chúng tôi kết luận rằng quý vị nhiễm COVID-19.

\item

  \textbf{Nếu xét nghiệm của quý vị cho biết ``Không Thể Xét Nghiệm''}, điều đó
  có nghĩa là không đủ vật liệu trong mẫu của quý vị để chúng tôi xét nghiệm
  hoặc đã xảy ra lỗi khi chúng tôi thực hiện xét nghiệm. Do chúng tôi không thể
  xét nghiệm mẫu của quý vị, chúng tôi không thể biết quý vị nhiễm hay không
  nhiễm vi-rút gây ra COVID-19.

\end{itemize}

\bigskip

\large \underline{\textbf{Các Nguồn Lực}}

Nếu quý vị muốn biết thêm thông tin về COVID-19, vui lòng xem các nguồn lực sau
từ các đối tác y tế công cộng:

\begin{itemize}
\item
  Public Health \textemdash Seattle \& King County (Cơ Quan Y Tế Cộng Đồng \textemdash Seattle \& Quận King): \\
  \url{https://www.kingcounty.gov/depts/health/covid-19/languages/vietnamese.aspx}

\item

  Washington State Department of Health (Sở Y Tế Tiểu Bang Washington):\\
  \url{https://www.doh.wa.gov/Emergencies/NovelCoronavirusOutbreak2020COVID19/vietnamese}

\item
  U.S. Centers for Disease Control and Prevention (CDC, Trung Tâm Kiểm Soát và
  Phòng Ngừa Dịch Bệnh Hoa Kỳ):\\
  \url{https://vietnamese.cdc.gov/coronavirus/2019-ncov/index.html}

\end{itemize}

Nếu quý vị muốn trao đổi với ai đó qua điện thoại, vui lòng gọi đến một trong
những số sau:

\begin{itemize}

\item

  \textbf{Các câu hỏi y tế liên quan đến COVID-19:} Hãy liên hệ với King County
  Call Center (Tổng Đài Quận King) từ 8 giờ sáng đến 7 giờ tối theo số
  \textbf{206-477-3977}

\item

  \textbf{Các câu hỏi không liên quan đến y tế về COVID-19, bao gồm các vấn đề
  về tuân thủ và liên quan đến doanh nghiệp:} Hãy liên hệ với King County
  COVID-19 Business and Community Information Line (Đường Dây Thông Tin Doanh
  Nghiệp và Cộng Đồng trong COVID-19 của Quận King), từ 8 giờ 30 phút sáng đến 4
  giờ 30 phút chiều, Thứ Hai - Thứ Sáu theo số \textbf{206-296-1608}

\item

  \textbf{Các câu hỏi chung về COVID-19 tại Tiểu Bang Washington:} Hãy liên hệ
  với Washington State Novel Coronavirus Call Center (Tổng Đài Vi-rút Corona
  Chủng Mới của Tiểu Bang Washington) từ 6 giờ sáng đến 10 giờ tối theo số
  \textbf{1-800-525-0127}

\end{itemize}

\bigskip

\large \underline{\textbf{Các Biện Pháp và Hạn Chế}}

\textbf{Biện Pháp}\\
Xét nghiệm này được thực hiện để phát hiện SAR-CoV-2 bằng cách sử dụng PCR phiên
mã ngược. Xét nghiệm này sử dụng kỹ thuật khuếch đại các phân đoạn vi-rút ORf1B
và S cũng như phát hiện kiểm soát khuếch đại RnaseP ở người. Northwest Genomics
Center của University of Washington đã phát triển và xác định các đặc tính hiệu
suất của loại xét nghiệm này. Phòng thí nghiệm này được WA Department of Health
Medical Testing Site (Sở Phụ Trách Cơ Sở Xét Nghiệm Y Tế Sức Khỏe) của WA cấp
phép là đủ tiêu chuẩn để thực hiện các xét nghiệm lâm sàng có độ phức tạp cao.

\textbf{Hạn Chế}

\begin{itemize}

\item

  Các kết quả này được cung cấp như một phần trong nghiên cứu.

\item

  Quý vị đã tự lấy mẫu và dán nhãn cho bản thân mình (hoặc làm cho ai đó), sau
  đó gửi lại cho phòng thí nghiệm qua dịch vụ giao hàng. Mẫu không được thu
  thập, dán nhãn hoặc giao nhận bởi một chuyên gia chăm sóc sức khỏe như y tá
  hoặc bác sĩ. Điều này có thể ảnh hưởng đến kết quả xét nghiệm của quý vị và
  một nhà cung cấp dịch vụ chăm sóc sức khỏe có thể muốn xét nghiệm lại cho quý
  vị.

\item

  Vì báo cáo này dựa trên dữ liệu từ một nghiên cứu, kết quả của quý vị có thể
  không chẩn đoán COVID-19, căn bệnh gây ra do nhiễm vi-rút SARS-CoV-2.

\end{itemize}

\end{document}
